Bayesian inference is rarely available in closed form, and the
standard computational response, Markov chain Monte Carlo, must typically be
re-run from scratch for every new dataset. Neural Bayes estimators (NBEs), a recent and fast-developing approach to simulation-based inference, 
avoid this per-dataset cost by amortisation with a neural network. A neural network is trained once, on simulated data, to approximate the Bayes estimator for some parameters under a Bayesian model, after which inference for any new dataset is a
single forward pass. This thesis begins from the observation that training such an estimator itself involves a Monte Carlo computation: the training objective is a Monte
Carlo estimate of the Bayes risk, and the gradient that drives
training is therefore a random quantity with a variance. Taking this view, we apply variance-reduction techniques to the training of the estimator to see whether they can reduce the computational burden of training.

Our first contribution is to show how Rao-Blackwellisation can be used to make training more efficient. Analogous to Gibbs sampling, this method is applicable when some conditional distributions of the full Bayesian joint distribution are available in closed form. When a block of the parameter vector has a tractable conditional posterior, that block is integrated out of the loss analytically rather than sampled; under squared-error loss this amounts to substituting the conditional posterior mean for the sampled target. The resulting training gradient is shown to be the Rao-Blackwellised variant of the standard one, so the law of total variance gives a positive semi-definite gradient-variance gap: the technique can never increase the variance of the training gradient, for any model or architecture. When applied to Bayesian linear regression as a case study, the Rao-Blackwellised estimator attains a lower loss than the standard NBE in every cell of an experimental grid, closes half of the gap to the Gibbs benchmark, and matches the standard NBE at a quarter of the minibatch size.

As a contrasting case, the thesis also studies importance sampling, which requires no tractable conditional posterior and instead changes the distribution of simulated training examples before reweighting them back to the original Bayes-risk objective. This greater generality comes without the same guarantee Rao-Blackwellisation had on gradient variance reduction. In an autoregressive experiment, importance sampling gives no accuracy improvement: the standard estimator is already close to the Bayes-optimal reference across most of the prior, and the proposals considered do not target the main source of remaining error. More aggressive reweighting instead lowers the effective sample size and increases noise. The results therefore separate two kinds of variance reduction for NBEs: Rao-Blackwellisation gives a guaranteed reduction in gradient variance when conjugacy is available, while importance sampling is more broadly applicable but problem-dependent.
